\paragraph{window-\texorpdfstring{$6$}{6} 微态偏置纤维律的 Cauchy--Stieltjes 闭式、Cartan--Lanczos 共轭与 Toda 轨道}
在 window-\(6\) 的三档纤维直方图、\(B_3/C_3\) 根字典、projective long sector 的纯 Cartan 晶场分解，以及极小紧 Lie 完成被唯一锁定到 \(\mathrm{Spin}(7)\) 之后，还可把同一离散对象进一步压缩为一条完全闭合的谱测度链：先由微态均匀分布推出纤维大小的偏置概率律，再由该三原子测度生成唯一的 Jacobi--Lanczos Hamiltonian，继而把它识别为同一 \(3\)-能级 Cartan 可观测量在特定量子态下的 Born 分布，最后再沿有限非周期 Toda 流把这一 Hamiltonian 放入同一 regular semisimple 轨道的等谱异宿链。由此，window-\(6\) 的直方图、Cartan 能级、Lanczos 链、正交多项式与 Toda 轨道不再是彼此并列的读法，而是同一谱测度对象的五种等价坐标。

\begin{theorem}[微态偏置纤维律的 Cauchy--Stieltjes/Herglotz 闭式]\label{thm:xi-time-part9zc-window6-microstate-biased-fiberlaw-herglotz-closedform}
沿用推论 \ref{cor:fold-bin-m6-fiber-hist} 的记号，定义 window-\(6\) 的微态偏置纤维律
\[
\nu_6
:=
\frac1{64}\sum_{w\in X_6}d_6(w)\,\delta_{d_6(w)}.
\]
则
\[
\nu_6=\frac14\,\delta_2+\frac3{16}\,\delta_3+\frac9{16}\,\delta_4.
\]
定义其 Cauchy--Stieltjes 变换
\[
G_6(z):=\int_{\RR}\frac{\dd\nu_6(t)}{z-t},
\qquad
M_6(z):=-G_6(z)=\int_{\RR}\frac{\dd\nu_6(t)}{t-z}.
\]
则 \(M_6\) 在上半平面上是 Herglotz 函数，而 \(G_6\) 具有精确闭式
\[
G_6(z)=\frac{16z^2-91z+126}{16(z-2)(z-3)(z-4)}.
\]
\leanverified{paper\_xi\_time\_part9zc\_window6\_microstate\_biased\_fiberlaw\_herglotz\_closedform}
等价地，若
\[
m_n:=\int_{\RR} t^n\,\dd\nu_6(t)
\qquad (n\ge 0),
\]
则
\[
m_n=\frac{8\cdot 2^{n+1}+4\cdot 3^{n+1}+9\cdot 4^{n+1}}{64},
\]
并且在 \(z=\infty\) 邻域内有 Laurent 展开
\[
G_6(z)=\sum_{n\ge 0}\frac{m_n}{z^{n+1}}.
\]
\end{theorem}
\begin{proof}
由推论 \ref{cor:fold-bin-m6-fiber-hist}，window-\(6\) 的纤维大小只取 \(2,3,4\) 三个值，且相应纤维个数分别为 \(8,4,9\)。总微态数为
\[
8\cdot 2+4\cdot 3+9\cdot 4=64,
\]
故微态均匀分布向纤维大小层的推前正是
\[
\nu_6(\{2\})=\frac{8\cdot 2}{64}=\frac14,\qquad
\nu_6(\{3\})=\frac{4\cdot 3}{64}=\frac3{16},\qquad
\nu_6(\{4\})=\frac{9\cdot 4}{64}=\frac9{16}.
\]
于是
\[
G_6(z)=\frac{1/4}{z-2}+\frac{3/16}{z-3}+\frac{9/16}{z-4}.
\]
将三项通分，得到
\[
G_6(z)=\frac{4(z-3)(z-4)+3(z-2)(z-4)+9(z-2)(z-3)}{16(z-2)(z-3)(z-4)}
=\frac{16z^2-91z+126}{16(z-2)(z-3)(z-4)}.
\]

另一方面，对任意 \(n\ge 0\)，
\[
m_n=\frac14\,2^n+\frac3{16}\,3^n+\frac9{16}\,4^n
=\frac{8\cdot 2^{n+1}+4\cdot 3^{n+1}+9\cdot 4^{n+1}}{64}.
\]
再由几何级数展开
\[
\frac1{z-t}=\sum_{n\ge 0}\frac{t^n}{z^{n+1}}
\qquad (|z|>|t|),
\]
逐项积分即得 \(G_6\) 在 \(z=\infty\) 邻域内的 Laurent 展开。最后，\(\Im z>0\) 时
\[
\Im M_6(z)=\int_{\RR}\frac{\Im z}{|t-z|^2}\,\dd\nu_6(t)>0,
\]
故 \(M_6\) 为 Herglotz 函数。证毕。
\end{proof}

\begin{theorem}[window-\texorpdfstring{$6$}{6} 微态偏置律的正交多项式闭合与三次终止]\label{thm:xi-time-part9zc-window6-orthogonal-polynomials-cubic-termination}
\leanverified{paper\_xi\_time\_part9zc\_window6\_orthogonal\_polynomials\_cubic\_termination}
以 \(\nu_6\) 为内积
\[
\langle f,g\rangle_{\nu_6}:=\int_{\RR}f(t)g(t)\,\dd\nu_6(t).
\]
令 \(P_0,P_1,P_2\) 为 \(\nu_6\) 的 monic 正交多项式，并令 \(P_3\) 为唯一与所有次数至多 \(2\) 的多项式正交的 monic 三次多项式。则
\[
P_0(x)=1,
\]
\[
P_1(x)=x-\frac{53}{16},
\]
\[
P_2(x)=x^2-\frac{371}{61}x+\frac{516}{61},
\]
\[
P_3(x)=(x-2)(x-3)(x-4).
\]
其非退化平方范数为
\[
h_0:=\langle P_0,P_0\rangle_{\nu_6}=1,\qquad
h_1:=\langle P_1,P_1\rangle_{\nu_6}=\frac{183}{256},
\qquad
h_2:=\langle P_2,P_2\rangle_{\nu_6}=\frac9{61},
\]
而
\[
\langle P_3,P_3\rangle_{\nu_6}=0.
\]
此外，\(P_2\) 的两个零点为
\[
x_{\pm}=\frac{371\pm 11\sqrt{97}}{122},
\]
并满足严格交错
\[
2<x_-<3<x_+<4.
\]
\end{theorem}
\begin{proof}
由定理 \ref{thm:xi-time-part9zc-window6-microstate-biased-fiberlaw-herglotz-closedform}，
\[
m_0=1,\qquad
m_1=\frac{53}{16},\qquad
m_2=\frac{187}{16},\qquad
m_3=\frac{689}{16}.
\]
故
\[
P_1(x)=x-m_1=x-\frac{53}{16},
\]
并且
\[
h_1=m_2-m_1^2=\frac{187}{16}-\left(\frac{53}{16}\right)^2=\frac{183}{256}.
\]

设
\[
P_2(x)=x^2+ux+v.
\]
由 \(\langle P_2,1\rangle_{\nu_6}=0\) 与 \(\langle P_2,x\rangle_{\nu_6}=0\)，得到线性方程组
\[
m_2+u m_1+v=0,
\qquad
m_3+u m_2+v m_1=0.
\]
代入上面的矩并求解，得到
\[
u=-\frac{371}{61},
\qquad
v=\frac{516}{61}.
\]
故
\[
P_2(x)=x^2-\frac{371}{61}x+\frac{516}{61}.
\]
再直接在支撑点 \(2,3,4\) 上代入，得到
\[
P_2(2)=\frac{18}{61},\qquad
P_2(3)=-\frac{48}{61},\qquad
P_2(4)=\frac8{61}.
\]
于是
\[
h_2
=
\frac14\left(\frac{18}{61}\right)^2
+
\frac3{16}\left(\frac{48}{61}\right)^2
+
\frac9{16}\left(\frac8{61}\right)^2
=\frac9{61}.
\]

由于 \(\operatorname{supp}\nu_6=\{2,3,4\}\) 恰含三点，唯一的 monic 三次多项式
\[
P_3(x):=(x-2)(x-3)(x-4)
\]
在支撑上处处为零，故自动与所有次数至多 \(2\) 的多项式正交，且其平方范数为零。

最后，对方程
\[
61x^2-371x+516=0
\]
求根即可得到
\[
x_{\pm}=\frac{371\pm 11\sqrt{97}}{122}.
\]
又因
\[
P_2(2)>0,\qquad P_2(3)<0,\qquad P_2(4)>0,
\]
故两根必分别落在区间 \((2,3)\) 与 \((3,4)\) 内。证毕。
\end{proof}

\begin{theorem}[window-\texorpdfstring{$6$}{6} 的规范信息 Hamiltonian]\label{thm:xi-time-part9zc-window6-canonical-information-hamiltonian}
定义三对角矩阵
\[
J_6:=
\begin{pmatrix}
\frac{53}{16} & \frac{\sqrt{183}}{16} & 0\\[1mm]
\frac{\sqrt{183}}{16} & \frac{2703}{976} & \frac{16\sqrt3}{61}\\[1mm]
0 & \frac{16\sqrt3}{61} & \frac{178}{61}
\end{pmatrix}.
\]
则：
\begin{enumerate}
  \item \(\chi_{J_6}(\lambda)=\det(\lambda I_3-J_6)=(\lambda-2)(\lambda-3)(\lambda-4)\)；
  \item \(\nu_6\) 正是 \(J_6\) 在循环向量 \(e_1=(1,0,0)^{\mathsf T}\) 下的谱测度；亦即对任意 Borel 集 \(B\subset\RR\),
  \[
  \nu_6(B)=\left\langle e_1,\mathbf 1_B(J_6)e_1\right\rangle;
  \]
  \item 等价地，对一切 \(n\ge 0\),
  \[
  \int_{\RR} t^n\,\dd\nu_6(t)=\left\langle e_1,J_6^n e_1\right\rangle.
  \]
\end{enumerate}
因此，window-\(6\) 的全部微态偏置信息可被唯一压缩为一个 \(3\)-站点 Jacobi Hamiltonian，而不损失任何矩数据。
\end{theorem}
\begin{proof}
由定理 \ref{thm:xi-time-part9zc-window6-orthogonal-polynomials-cubic-termination}，
\[
P_1(x)=\left(x-\frac{53}{16}\right)P_0(x),
\]
\[
P_2(x)=\left(x-\frac{2703}{976}\right)P_1(x)-\frac{183}{256}P_0(x),
\]
\[
P_3(x)=\left(x-\frac{178}{61}\right)P_2(x)-\frac{768}{3721}P_1(x).
\]
这三式分别把三项递推系数读出为
\[
a_0=\frac{53}{16},\qquad
\beta_1=\frac{183}{256},\qquad
a_1=\frac{2703}{976},\qquad
\beta_2=\frac{768}{3721},\qquad
a_2=\frac{178}{61}.
\]
于是对应的 Jacobi 矩阵恰为
\[
J_6=
\begin{pmatrix}
a_0 & \sqrt{\beta_1} & 0\\
\sqrt{\beta_1} & a_1 & \sqrt{\beta_2}\\
0 & \sqrt{\beta_2} & a_2
\end{pmatrix},
\]
即题设矩阵。

对有限支撑正测度的 Favard 理论表明，在正交归一基下，“乘以 \(x\)” 的算子由该 Jacobi 矩阵表示，因此其特征多项式正是递推终止多项式 \(P_3\)。由定理 \ref{thm:xi-time-part9zc-window6-orthogonal-polynomials-cubic-termination}，
\[
P_3(\lambda)=(\lambda-2)(\lambda-3)(\lambda-4),
\]
从而第一条成立。

同一谱定理同时给出
\[
\left\langle e_1,J_6^n e_1\right\rangle
=
\int_{\RR}t^n\,\dd\nu_6(t)
\qquad (n\ge 0),
\]
于是 \(\nu_6\) 正是 \(J_6\) 在 \(e_1\) 下的谱测度。证毕。
\end{proof}

\begin{theorem}[spin-\texorpdfstring{$1$}{1} Cartan 可观测量的 Born 实现]\label{thm:xi-time-part9zc-window6-spin1-cartan-born-realization}
沿用定理 \ref{thm:conclusion-window6-projective-long-sector-spin1-crystal-field} 在 \(q=1\) 时的记号，定义
\[
C_6:=3I_3+J_z=\operatorname{diag}(2,3,4),
\qquad
J_z=\operatorname{diag}(-1,0,1).
\]
再定义单位向量
\[
\psi_6:=
\begin{pmatrix}
\frac12\\[1mm]
\frac{\sqrt3}{4}\\[1mm]
\frac34
\end{pmatrix},
\qquad
\|\psi_6\|=1.
\]
则 \(\nu_6\) 恰为 \(C_6\) 在态 \(\psi_6\) 下的谱律：
\[
\nu_6(B)=\left\langle \psi_6,\mathbf 1_B(C_6)\psi_6\right\rangle
\qquad (B\subset\RR\ \text{Borel}),
\]
亦即
\[
\int_{\RR} t^n\,\dd\nu_6(t)=\left\langle \psi_6,C_6^n\psi_6\right\rangle
\qquad (n\ge 0).
\]
因此，window-\(6\) 的微态偏置纤维律严格等于一个 spin-\(1\) 系统对 Cartan 可观测量 \(3I_3+J_z\) 的 Born 测量分布。
\end{theorem}
\begin{proof}
由于 \(C_6\) 已在标准基下对角化，其三个谱投影分别投到三条坐标轴。因此
\[
\left\langle \psi_6,\mathbf 1_{\{2\}}(C_6)\psi_6\right\rangle
=
\left|\frac12\right|^2
=
\frac14,
\]
\[
\left\langle \psi_6,\mathbf 1_{\{3\}}(C_6)\psi_6\right\rangle
=
\left|\frac{\sqrt3}{4}\right|^2
=
\frac3{16},
\]
\[
\left\langle \psi_6,\mathbf 1_{\{4\}}(C_6)\psi_6\right\rangle
=
\left|\frac34\right|^2
=
\frac9{16}.
\]
这恰与定理 \ref{thm:xi-time-part9zc-window6-microstate-biased-fiberlaw-herglotz-closedform} 中的 \(\nu_6\) 一致。矩恒等式再由谱定理立即推出。证毕。
\end{proof}

\begin{theorem}[Cartan--Lanczos 共轭定理]\label{thm:xi-time-part9zc-window6-cartan-lanczos-conjugacy}
\leanverified{paper\_xi\_time\_part9zc\_window6\_cartan\_lanczos\_conjugacy}
存在唯一正交矩阵 \(U\in SO(3)\)，其第一行等于 \(\psi_6^{\mathsf T}\)，并满足
\[
J_6=U\,C_6\,U^{\mathsf T}.
\]
更具体地，可取
\[
U=
\begin{pmatrix}
\frac12 & \frac{\sqrt3}{4} & \frac34\\[1mm]
-\frac{7\sqrt{183}}{122} & -\frac{5\sqrt{61}}{244} & \frac{11\sqrt{183}}{244}\\[1mm]
\frac{3\sqrt{61}}{61} & -\frac{4\sqrt{183}}{61} & \frac{2\sqrt{61}}{61}
\end{pmatrix}\in SO(3),
\]
从而
\[
J_6=U\,\operatorname{diag}(2,3,4)\,U^{\mathsf T}.
\]
因此，定理 \ref{thm:conclusion-window6-projective-long-sector-spin1-crystal-field} 中的 Cartan 势与定理 \ref{thm:xi-time-part9zc-window6-canonical-information-hamiltonian} 中的 Jacobi Hamiltonian 并非两件不同对象，而是同一 regular semisimple 轨道在 Cartan 基与 Krylov--Lanczos 基中的两种坐标表达。
\end{theorem}
\begin{proof}
令
\[
\phi_0(x)=1,
\qquad
\phi_1(x)=\frac{P_1(x)}{\sqrt{h_1}},
\qquad
\phi_2(x)=\frac{P_2(x)}{\sqrt{h_2}},
\]
其中 \(P_k,h_k\) 见定理 \ref{thm:xi-time-part9zc-window6-orthogonal-polynomials-cubic-termination}。再记
\[
(\lambda_1,\lambda_2,\lambda_3)=(2,3,4),
\qquad
(w_1,w_2,w_3)=\left(\frac14,\frac3{16},\frac9{16}\right).
\]
定义
\[
U_{ij}:=\sqrt{w_j}\,\phi_{i-1}(\lambda_j)
\qquad (1\le i,j\le 3).
\]
由于 \((\phi_0,\phi_1,\phi_2)\) 在 \(\nu_6\) 下正交归一，故 \(U\) 为正交矩阵。其第一行为
\[
(\sqrt{w_1},\sqrt{w_2},\sqrt{w_3})
=
\left(\frac12,\frac{\sqrt3}{4},\frac34\right)
=
\psi_6^{\mathsf T}.
\]

另一方面，在正交多项式基下，“乘以 \(x\)” 的矩阵表示正是 \(J_6\)，而在谱基下则是 \(\operatorname{diag}(2,3,4)\)。因此
\[
J_6=U\,\operatorname{diag}(2,3,4)\,U^{\mathsf T}.
\]
将 \(\phi_1,\phi_2\) 在三点 \(2,3,4\) 上的取值显式代入，便得到题设矩阵。再直接计算可得 \(\det U=1\)，故 \(U\in SO(3)\)。

最后，由于 \(J_6\) 的谱简单，且各列是按特征值 \(2,3,4\) 排列并要求第一行全部取正，因此每一列的符号都已被唯一固定，从而满足题设条件的正交矩阵只能有这一组。证毕。
\end{proof}

\begin{corollary}[三层自能的精确 \texorpdfstring{$J$}{J}-fraction]\label{cor:xi-time-part9zc-window6-three-layer-selfenergy-jfraction}
定理 \ref{thm:xi-time-part9zc-window6-microstate-biased-fiberlaw-herglotz-closedform} 中的 Cauchy--Stieltjes 变换 \(G_6\) 具有精确 Jacobi continued fraction
\[
G_6(z)
=
\frac{1}{
z-\frac{53}{16}
-\cfrac{\frac{183}{256}}{
z-\frac{2703}{976}
-\cfrac{\frac{768}{3721}}{
z-\frac{178}{61}
}}}.
\]
等价地，若定义
\[
\Sigma_2(z):=\frac{\frac{768}{3721}}{z-\frac{178}{61}},
\qquad
\Sigma_1(z):=\frac{\frac{183}{256}}{z-\frac{2703}{976}-\Sigma_2(z)},
\]
则
\[
G_6(z)=\frac{1}{z-\frac{53}{16}-\Sigma_1(z)}.
\]
因此，window-\(6\) 的微态偏置纤维律可被解释为一个严格三层终止的自能递归。
\end{corollary}
\begin{proof}
由定理 \ref{thm:xi-time-part9zc-window6-canonical-information-hamiltonian},
\[
G_6(z)=\left\langle e_1,(zI_3-J_6)^{-1}e_1\right\rangle.
\]
对三对角矩阵 \(zI_3-J_6\) 作两次 Schur 补，得到
\[
G_6(z)
=
\frac{1}{
z-a_0-\cfrac{\beta_1}{z-a_1-\cfrac{\beta_2}{z-a_2}}},
\]
其中
\[
a_0=\frac{53}{16},\qquad
a_1=\frac{2703}{976},\qquad
a_2=\frac{178}{61},\qquad
\beta_1=\frac{183}{256},\qquad
\beta_2=\frac{768}{3721}.
\]
代入即可得到所示 continued fraction。再把它化简，便回到定理 \ref{thm:xi-time-part9zc-window6-microstate-biased-fiberlaw-herglotz-closedform} 的有理闭式。证毕。
\end{proof}
\endinput
